We opened with the observation that models achieving 95\%+ benchmark accuracy can fail across multiple trustworthiness dimensions simultaneously at near-identical rates ($r > 0.99$)---a model robust to adversarial attacks failing reliability and fairness tests proportionally. Our large-scale correlation analysis across 20 LLMs confirms this observed coupling ($r > 0.99$, $p < 2 \times 10^{-17}$) is robust within our 3-benchmark sample, persistent across model scales ($r > 0.98$ within size strata), and far stronger than the moderate correlations ($r > 0.3$) expected from shared root causes. This finding challenges the independent dimensions assumption underlying current trustworthiness evaluation, where reliability (TrustfulQA), robustness (AdvBench), and fairness (BOLD) are assessed separately without testing empirical orthogonality. However, interpretation remains uncertain---correlation cannot distinguish whether benchmarks measure unified constructs or distinct dimensions with shared confounds.

However, our attempt to map this coupling into a validated failure mode taxonomy failed: hierarchical clustering produced perfectly stable groups (100\% bootstrap consistency) but poor separation (silhouette = 0.274 $<$ 0.5), revealing that 3-benchmark evaluation cannot resolve distinct failure modes despite reproducible structure. The clustering refutation is not a flaw in methodology but a hard constraint---testing $k \geq 3$ failure modes requires $n \geq 3$ benchmarks, yet our $k=2$ evaluation (TrustfulQA+AdvBench vs BOLD) lacks both sample size and separation quality for taxonomy validation. This negative result is methodologically informative: it establishes that correlation analysis and clustering analysis have different sample size requirements, with $n=20$ models sufficient for correlation statistical power but $n=3$ benchmarks insufficient for clustering beyond binary partitions.

The near-perfect correlations ($r > 0.99$) admit two competing explanations, each with distinct implications for future work. The \textbf{Unified Capability Hypothesis} (appears more plausible given alignment with prior RLHF observations, but not quantifiable without Bayesian analysis) posits that reliability, robustness, and fairness---as currently operationalized---measure the same underlying construct rather than orthogonal dimensions. If validated through broader measurement, dimension-independent evaluation would be fundamentally misspecified, and practitioners should focus on interventions improving general model quality rather than dimension-specific fixes---though intervention validation remains untested (P4). The \textbf{Insufficient Resolution Hypothesis} (cannot be ruled out with current 3-benchmark data) posits that distinct failure modes exist but require broader coverage (5-10 dimensions) to reveal structure invisible in our limited sample. If validated, HELM's 50+ benchmark aggregation should shift toward correlation-aware evaluation, identifying which dimensions cluster together and which remain orthogonal.

These hypotheses make opposing predictions testable via concrete future work:

\textbf{FW1 (High Priority): Expand to 5-10 Benchmarks.} Add ToxiGen (toxicity), BBQ (bias), HellaSwag (reasoning), GSM8K (math), HumanEval (code) to test discriminating predictions: Unified Capability predicts $r > 0.95$ for $>80\%$ of pairs; Insufficient Resolution predicts mean $r < 0.85$ with $\geq 30\%$ of pairs $< 0.7$. This directly addresses the sample size limitation (L1) and enables $k=3-5$ clustering evaluation. Gates all other future work.

\textbf{FW4 (High Priority): Test Unified Capability via Factor Analysis.} Apply PCA or factor analysis to 10-benchmark data. If first component explains $>90\%$ variance, unified construct confirmed. If 3+ factors needed for 70\% variance, distinct dimensions emerge.

\textbf{FW5 (Medium Priority): Scale Invariance Direct Test.} Apply Mantel test to existing stratified correlation matrices (already computed in Section 5.1) to validate whether $r > 0.98$ within-strata correlations translate to Mantel $r > 0.7$ between-strata structure similarity.

\textbf{FW2 (Medium Priority): Revise Clustering Metrics.} Test alternative cluster quality metrics (Calinski-Harabasz, Davies-Bouldin) and relaxed thresholds (silhouette $> 0.2$ for $r > 0.95$ data) to determine whether bootstrap-silhouette divergence reflects metric inappropriateness or true absence of separation.

\textbf{FW6 (Low Priority): Intervention Validation on Highest-Correlation Pair.} Test whether calibration training targeting TrustfulQA (reliability) improves AdvBench (robustness) given $r = 0.998$ coupling. If improvement transfers (Cohen's $d > 0.5$ for both), validates shared root cause mechanism even without taxonomy.

The path forward is empirically clear: expand to 5-10 benchmarks to test whether trustworthiness is fundamentally unidimensional ($r > 0.99$ persists, factor analysis yields single dominant component) or our measurement was too coarse (correlations drop, distinct clusters emerge with silhouette $> 0.5$). \textbf{Implications:} If trustworthiness is unidimensional, current practice wastes resources evaluating 50+ near-redundant benchmarks---effort should shift toward broader construct coverage. If dimensions exist but are unresolved, benchmarks must expand to $\geq 10$ metrics minimum to distinguish them. Until this question resolves, practitioners should interpret multi-dimensional evaluation cautiously---correlation structure matters, not just aggregate scores.

\textbf{Closing the loop:} We opened with a puzzle (why do models fail multiple dimensions simultaneously?) and hypothesized it reflected distinct failure modes sharing root causes. The data reveal tighter coupling than hypothesized ($r > 0.99$ instead of $r > 0.3$) but refute the taxonomy claim (silhouette = 0.274 despite 100\% stability). The observed coupling is robust within our 3-benchmark sample, but interpretation remains uncertain pending broader measurement. The field faces a more fundamental question than we initially posed: are reliability, robustness, and fairness separable constructs, or do current benchmarks measure facets of a unified dimension? Our 3-benchmark design cannot answer this---but a 10-benchmark follow-up can.
